\subsection{Ablation outcomes}
\label{sec:ablationoutcomes}

\textbf{Every level of the sweep is judged by where its attempts end, not by
whether they return an image.} The sweep submits each level of each axis with
all three body configurations, and the gateway keeps a record of every task it
was handed, including the ones whose replies never arrived.
Figure~\ref{fig:sweep} reports that record. Three outcomes are distinguished:
the request is declined outright, the request succeeds but the generated image
is withheld, and the attempt does not conclude at all.

\begin{figure}[t]
\centering
\includegraphics[width=\linewidth]{figures/fig_sweep_45.png}
\caption{\textbf{Where each attempt of the ablation sweep ended.} Each row is
one level of one axis. Segments count attempts, and the right-hand figure gives
images returned over resolved attempts at that level. Grey attempts have not
concluded and are not counted as refusals; a level with a wide grey band has
mostly not been measured, which the figure is drawn to make visible rather than
absorb into a rate.}
\label{fig:sweep}
\end{figure}

The two refusal kinds are not interchangeable. A declined request is stopped
before a picture exists, so its rate says nothing about the requested
composition. A withheld output was drawn and then rejected, so its rate is
evidence that the composition itself carried the refusal. Counting both into a
single refusal rate would merge two different findings, and counting only
returned images discards the more informative of the two.

\textbf{The column that separates the axes from the comparison arms is the
declined-request column.} The axis levels decline few requests and return
almost no images; the comparison arms decline many, and they account for every
image the sweep returns. This is the pattern the method section predicts: the
axes change how the requested figure is dressed, which leaves the request still
reading as a request for the same picture, whereas the comparison arms change
what is being asked for. Re-dressing the composition does not alter whether the
prompt layer stops the request, and altering that does not by itself return a
picture.

\begin{table}[t]
\centering\small
\setlength{\tabcolsep}{5pt}
\renewcommand{\arraystretch}{1.12}
\begin{tabularx}{\linewidth}{@{}Xrrr@{}}
\toprule
Arm & Images & Declined & Withheld \\
\midrule
Axis levels (seven axes) & 9/83 & 1 & 68 \\
Comparison arms (five) & 4/45 & 12 & 27 \\
\midrule
Fisher exact, declined requests & \multicolumn{3}{r}{$p < 0.0001$} \\
Fisher exact, images returned & \multicolumn{3}{r}{$p = 1.000$} \\
\bottomrule
\end{tabularx}
\caption{\textbf{The two arms differ in where their refusals happen, not in
how often a picture comes back.} Denominators are attempts. The arms are
indistinguishable on image return and separated by a wide margin on declined
requests.}
\label{tab:arms}
\end{table}

\textbf{The comparison arms are counted on the same footing as the axes.}
Against the four prior-work arms, the method of this paper returns four images
in twelve attempts while the four return none in forty-eight; the pooled
comparison is in Table~\ref{tab:arms} and the four are counted individually in
Appendix~\ref{app:ablationdetail}. The difference is large enough to read from
these counts alone, and the unit is the one used everywhere else in this paper:
one attempt is one submission.

\subsection{Why elapsed time is not used}
\label{sec:notiming}

An earlier version of this analysis classified refusals by how long they took: a
refusal inside ninety seconds was read as a request stopped before generation,
and a slower one as a generated image withheld. That rule is withdrawn, and
Figure~\ref{fig:rules} shows why.

\begin{figure}[t]
\centering
\includegraphics[width=0.86\linewidth]{figures/fig_rules_45.png}
\caption{\textbf{The same attempts under two classification rules.} The
retired rule moves forty-three attempts from one group into the other. Elapsed
time measures how long a request waited before an account picked it up, which
depends on the queue and on the account pool rather than on the prompt.}
\label{fig:rules}
\end{figure}

The rule fails on its own terms. The same cell, with the same refusal text, was
once answered in 54 seconds and once in 1896. Across cells, the ratio between a
level's fastest and slowest refusal reaches two orders of magnitude. Timing
therefore does not separate the two refusal kinds, and any rate computed on it
describes the queue rather than the content. The classification used here reads
the refusal text, and the two phrasings are distinct: one declines the request,
the other objects to the image that was produced.

\textbf{No number in this paper depends on timing.} Returned images are counted
as images; refusals are counted by their wording; attempts that reach no outcome
are reported separately and excluded from every denominator. The gateways lose
requests at very different rates, which is why the denominator is attempts
rather than resolved judgements: over attempts the two gateways agree, and over
judgements they differ by an order of magnitude.

\textbf{What the sweep does and does not establish.} It establishes that the
axis levels behave alike on the outcome that matters, so the axes do not rank
against one another, and that the comparison arms differ from the axes in where
their refusals occur. It does not establish an ordering among the axis levels.
At the coverage reached so far, the per-level counts are too small to separate a
level that returns pictures at an elevated rate from a baseline that returns
few: a two-sided test over the resolved levels finds no level distinguishable
from the pooled baseline, and the counts needed to change that are several times
what has been collected.

